\chapter{Introduction}

Cyber Security is about keeping computing systems working as intended, keeping data accessed only as desired, and keeping the computing resources secured. All of these need to be done in the presence of an adversary and on a budget. We first introduce some concepts in this chapter. 

\textbf{Threat Model}

A threat is a potential violation of security by an attacker. Since it is not possible to protect against all possible attackers, we need risk management. Computer security is always related to a specific threat model. 

Threat modelling works to identify, communicate and understand threats and mitigations within the context of protecting something of value. A threat model includes:

- What to protect

- From whom to protect

- Assumptions (can be wrong)

- A list of the (most important) relevant threats

- What actions are to be taken for each threat

\textbf{Security Policy} 

This is a statement on what is and what is not allowed. 

\textbf{Security Mechanism}

This is a method, tool or procedure for enforcing a security policy. A security policy can be implemented by a variety of security mechanisms. For a security mechanism to work, 

1. the security policy correctly describes the security requirement of a system

2. the security mechanism is implemented correctly

3. the security mechanism is deployed correctly

4. the security mechanism cannot be bypassed 

\textbf{Prevention} 

Prevention aims to prevent something bad from happening. 

\textbf{Detection} 

Detection is about knowing that something is wrong. There are two special cases, which are false positive, meaning this is a wrong alarm, and FALSE NEGATIVE, meaning there is an unnoticed attack.

\textbf{Response} 

Response is needed when we detect something is going wrong. 

\textbf{Mitigation} 

Mitigation can be more economical and practical than prevention. 